% BHU PYQ -- Professional Exam Template
\documentclass[11pt,a4paper]{article}

\usepackage[a4paper,top=2.5cm,bottom=2.5cm,left=2.8cm,right=2.8cm,headheight=14pt]{geometry}
\usepackage{amsmath,amssymb}
\usepackage[T1]{fontenc}
\usepackage[utf8]{inputenc}
\usepackage{lmodern}
\usepackage{microtype}
\usepackage{fancyhdr}
\usepackage{enumitem}
\usepackage{hyperref}
\usepackage{lastpage}
\usepackage{parskip}
\usepackage{tikz}
\usetikzlibrary{arrows.meta,shapes,positioning}

\hypersetup{colorlinks=true,urlcolor=black,linkcolor=black,
  pdftitle={B.Sc. (Hons.) Zoology Sem-V ZOB-502 -- BHU PYQ}}

\pagestyle{fancy}
\fancyhf{}
\renewcommand{\headrulewidth}{0.4pt}
\renewcommand{\footrulewidth}{0.4pt}
\fancyhead[L]{\small   \textit{Banaras Hindu University}}
\fancyhead[R]{\small   \textit{Zoology Paper: ZOB-502}}
\fancyfoot[C]{\small Page  \thepage\ of \pageref{LastPage}}


\newlist{parts}{enumerate}{3}
\setlist[parts,1]{label=(\alph*),leftmargin=*,itemsep=4pt,topsep=3pt}
\setlist[parts,2]{label=(\roman*),leftmargin=*,itemsep=2pt,topsep=2pt}
\setlist[parts,3]{label=(\arabic*),leftmargin=*,itemsep=2pt,topsep=2pt}

\newcommand{\pts}[1]{\hfill   \textit{[#1]}}


\begin{document}


\begin{center}
  {\large\bfseries BANARAS HINDU UNIVERSITY}\\[2pt]
  {\normalsize B.Sc. (Hons.) Semester V Examination, 2013-14}\\[6pt]
  \rule{0.8\linewidth}{0.5pt}\\[6pt]
  {\large\bfseries Zoology}\\[4pt]
  {\large\bfseries Paper: ZOB-502 --- Biochemistry \& Mammalian Physiology}\\[4pt]
  {\normalsize Time: Three hours \hfill Full Marks: 70}
\end{center}

\rule{\linewidth}{0.4pt}

\smallskip



\noindent   \textit{Note: Answer three questions from Section-A and three questions from Section-B. Question No. 1 in both sections is compulsory. Use a separate answer book for each Section. The figures in the right-hand margin indicate marks.}

\bigskip

\begin{center}
   \textbf{SECTION -- A} \\
   \textbf{(Biochemistry) (Marks: 35)}
\end{center}

\smallskip



\noindent   \textit{Note: Answer three questions including Question No. 1 which is compulsory.}

\medskip


\noindent   \textbf{Question 1.} Answer the following parts:

\begin{parts}
  \item Indicate whether the following statements are True or False. Justify your answer precisely:\pts{$4  \times 2 = 8$}
  
\begin{parts}
    \item Formation of ES complex is a compulsory step of enzyme catalysis.
    \item Cholesterol is composed of more than one fatty acids.
    \item One DNA Pol-III complex catalyzes DNA synthesis simultaneously at both strands of the parental DNA.
    \item `Shine Dalgarno' sequences determine initiation of prokaryotic transcription.
  \end{parts}

  \item Define the following:\pts{$3  \times 1 = 3$}
  
\begin{parts}
    \item Genetic Vector
    \item pI of an amino acid
    \item Acid-base catalysis.
  \end{parts}
\end{parts}

\medskip


\noindent   \textbf{Question 2.} Answer the following parts:\pts{$6 + 6 = 12$}

\begin{parts}
  \item Illustrate stepwise formation of transcription initiation complex for RNA pol II activity.
  \item Describe (in diagrammatic form only) four step events of chain elongation during prokaryotic protein synthesis.
\end{parts}

\medskip


\noindent   \textbf{Question 3.} Answer the following parts:\pts{$4 + 6 + 2 = 12$}

\begin{parts}
  \item Write the reactions producing FADH$_2$ and GTP during the TCA cycle.
  \item Illustrate arrangement of ETS at inner mitochondrial membrane.
  \item Construct LB-plot of the enzyme kinetics.
\end{parts}

\medskip


\noindent   \textbf{Question 4.} Write short notes on the following:\pts{$3  \times 4 = 12$}

\begin{parts}
  \item Prostaglandins
  \item tRNA
  \item Restriction endonucleases.
\end{parts}

\pagebreak

\begin{center}
   \textbf{SECTION -- B} \\
   \textbf{(Mammalian Physiology) (Marks: 35)}
\end{center}

\smallskip



\noindent   \textit{Note: Answer three questions including Question No. 1 which is compulsory.}

\medskip


\noindent   \textbf{Question 1.} Answer the following parts:

\begin{parts}
  \item Select and write the most suitable answer for the following:\pts{$2  \times \frac{1}{2} = 1$}
  
\begin{parts}
    \item Muscle contraction consists of the cyclic attachment and detachment of the:
    
\begin{parts}
      \item globular head portion of myosin to the F-actin filament
      \item tropomyosin to G-actin
      \item myosin head to G-actin
      \item ATP to actin-myosin complex
    \end{parts}

    \item 2,3-diphosphoglycerate is a charged anion which binds to:
    
\begin{parts}
      \item $\alpha$-chains of deoxyhaemoglobin
      \item $\beta$-chains of deoxyhaemoglobin
      \item $\alpha$-chains of oxy-haemoglobin
      \item $\beta$-chains of oxy-haemoglobin
    \end{parts}
  \end{parts}

  \item State whether the following statements are True or False. Write brief justifications for your answers:\pts{$2  \times 3 = 6$}
  
\begin{parts}
    \item Walls of afferent arteries contain rennin-secreting juxtaglomerular cells.
    \item Between trachea and alveolar sacs, airways divide 20 times and the first 16 generations of passages form the conducting zone of the airways.
    \item In human, Hering-Breuer reflex is not activated until tidal volume increases to greater than around 1.5 L.
  \end{parts}

  \item What is the speed of conduction if the distance between cathodal stimulating electrode and exterior electrode is 4 cm and the latent period is 2 ms?\pts{2}

  \item Define the following:\pts{2}
  
\begin{parts}
    \item Haemostasis
    \item GFR.
  \end{parts}
\end{parts}

\medskip


\noindent   \textbf{Question 2.} Answer the following parts:\pts{$6 + 6 = 12$}

\begin{parts}
  \item Describe the role of Ca$^{2+}$ during muscle contraction with suitable diagrams.
  \item Describe the role of pancreatic juice in digestion of proteins and carbohydrates.
\end{parts}

\medskip


\noindent   \textbf{Question 3.} Answer the following parts:\pts{$6 + 6 = 12$}

\begin{parts}
  \item Describe the role of respiratory muscles in breathing.
  \item What are the different types of synapse? Describe the process of nerve impulse propagation in a myelinated nerve-fiber.
\end{parts}

\medskip


\noindent   \textbf{Question 4.} Write short notes on the following:\pts{$3  \times 4 = 12$}

\begin{parts}
  \item Action potential
  \item Sarcomere
  \item Chloride shift.
\end{parts}

\vfill
\rule{\linewidth}{0.4pt}

\begin{center}\small
\href{https://bhu-exam.vercel.app/}{Click Here}\\[4pt]\href{https://me-aryan.vercel.app/}{Click Here}\\[4pt]\href{https://quashan-me.vercel.app/}{Click Here}
\end{center}

\end{document}
